% The DC/GTC report format is particularly created for M.S. and Ph.D. students at IITM. The idea is to encourage students to use Overleaf or other Latex compilers for creating a report. Since most of the students are unexposed to Latex, this will be an easy start to beginners which will eventually lead them to write their thesis, proposals and articles using Overleaf or Latex.

% Please don't get bogged down by the commands used. These are simple commands and can be used for multiple reports. To start with please concentrate on filling the blank spaces or editing the current text. Once you get a feel for using the report, you can start modifying the commands if required. 

% Please feel free to add or modify if required. If you have queries feel free to contact me at rameshkk@iitm.ac.in (Ramesh Kannan K)

\documentclass[12pt]{article}
\usepackage{times}
\usepackage{babel}
\usepackage{titling}
\usepackage{blindtext}
\usepackage{amsmath}
\usepackage{graphicx}
\usepackage{setspace}
\usepackage{natbib,etoolbox,lipsum,hyperref}


\usepackage[%
    left=2cm,%
    right=2cm,%
    top=3cm,%
    bottom=3cm,%
    paperheight=29.7cm,%
    paperwidth=21cm%
]{geometry} % Page settings

\begin{document}
\begin{titlepage}


\setlength{\droptitle}{-4em} 
\center

{\huge \bfseries Title of the proposed project..............} % Title of the project
\\[2cm] 


\begin{minipage}{0.4\textwidth}
\begin{center} \large
Submitted by: \textbf{NSS} % Name of the student
\\[0.5cm]
\textbf{19210051} % Registration number of the student
\end{center}
\end{minipage} 

%\addtolength{\droptitle}{-4pt}
\vspace{10em}

\includegraphics[width=6cm]{Logo.png}

\vspace{10em}
\begin{minipage}{0.6\textwidth}
\begin{center} \large
\textbf{Centre for Cognitive and Brain Sciences} % Name of the department
\\[0.5cm]
\textbf{Independent Seminar Report} % Type of meeting
\end{center}
\end{minipage} 

\vspace{5em}
{\large \today}


\end{titlepage}



\newpage
\begin{center}
\title*{\LARGE \textbf{Title of the proposed project}}       
\end{center}



\begin{flushleft}
\section*{\large {Introduction}}
\end{flushleft}


\newpage
\begin{flushleft}
\section*{\large {Literature review}}
The studies carried out by \cite{Smith:2012} and \citep{Smith:2013}

\end{flushleft}

\begin{figure}[!htbp]
    \centering
    \includegraphics[width=10cm]{image_1.jpg}
    \caption{Caption}
    \label{fig:my_label}
\end{figure}



\newpage
\begin{flushleft}
\section*{\large {Objectives}}


The objective of the proposed project are, 
\begin{itemize}
    \item First objective
    \item Second objective
\end{itemize}
\end{flushleft}

\newpage
\begin{flushleft}
\section*{\large {Methodology}}
\end{flushleft}

\newpage
\renewcommand\refname{\vskip -1cm}
\begin{flushleft}
\section*{\large {References}}
\bibliographystyle{harvard}
\bibliography{sample.bib}
\end{flushleft}

% We thank Prof. Allan A. Struthers at MTU for the style file for the references
\end{document}
